\chap{Changing what is prescribed}\label{sec:prescription}
The observable is part of a realization problem, not an annotation added after it has been solved. Prescribing more views can remove ambiguity, but it can also exclude every completion. This section follows both effects through the actual obstruction groups. It then works through a three-view model in which an initially harmless choice creates a genuine dead end.

\subsection{Three maps that must not be confused}
Keep the additive view system of Section~\ref{sec:obstruction}. For a prescribed domain \(U\), there are three distinct maps:
\[
 p_U:E\longrightarrow\prod_{u\in U}A_u,\qquad
 r_U:\Gamma(A)\longrightarrow\Gamma_U(A),\qquad
 \mathfrak o_U:\Gamma_U(A)\longrightarrow Q_U.
\]
The first restricts a raw family, whether or not it is coherent. The second restricts a global coherent family. The third detects whether an already coherent partial family has a global completion. Their domains are different. In particular, surjectivity of the raw coordinate projection does not imply surjectivity of \(r_U\).

The obstruction group is \(Q_U=R/B_U\), where \(B_U=\delta(\ker p_U)\). The ambiguity group is \(H_U=\ker\delta\cap\ker p_U\). Both depend on what has been fixed. Enlarging the prescription reduces the permitted correction space and the invisible exact space at the same time.

\subsection{Forgetting a prescription gives a quotient map}
Let \(U\subseteq W\). Then
\[
 \ker p_W\subseteq\ker p_U,\qquad B_W\subseteq B_U,
 \qquad H_W\subseteq H_U.
\]
A correction invisible on all of \(W\) is certainly invisible on \(U\). The direction of the quotient map is consequently
\[
 \pi_{W,U}:Q_W\longrightarrow Q_U,\qquad
 [r]_{B_W}\longmapsto[r]_{B_U}.
\]
The larger prescribed domain appears in the source. More protection means a smaller subgroup is divided out, so more obstruction information can remain.

\begin{theorem}[Functorial forgetting of protected views]\label{KS59}
The homomorphism \(\pi_{W,U}\) is well defined and surjective. For \(U\subseteq W\subseteq Z\), it satisfies
\[
 \pi_{U,U}=I,\qquad \pi_{W,U}\pi_{Z,W}=\pi_{Z,U}.
\]
For every \(s\in\Gamma_W(A)\),
\begin{equation}\label{eq:prescribed-naturality}
 \pi_{W,U}\mathfrak o_W(s)=\mathfrak o_U(s|_U).
\end{equation}
In particular an obstructed restriction remains obstructed when more compatible prescribed data are retained. A successful extension on the larger prescribed domain gives a successful extension of its restriction.
\end{theorem}
\begin{proof}
If \(r-r'\in B_W\), then \(r-r'\in B_U\), proving well-definedness. Every class modulo \(B_U\) has a representative \(r\), and the class of that same representative modulo \(B_W\) maps to it. This proves surjectivity. Identity and composition follow by applying the formulas to a representative.

Choose a raw extension \(\widetilde s\) of the partial section on \(W\). It is also a raw extension of \(s|_U\). Lift independence of the two extension obstructions therefore identifies both sides of \eqref{eq:prescribed-naturality} with the class of \(\delta\widetilde s\) modulo \(B_U\). A homomorphism sends zero to zero. Thus nonvanishing on the smaller domain rules out vanishing on the larger one, while vanishing on the larger domain implies vanishing after restriction.
\end{proof}
\source{KS59}
The converse implication does not hold. The quotient map can annihilate the obstruction created by a newly fixed value. Surjectivity of \(\pi_{W,U}\) is also not a theorem that arbitrary larger partial sections extend. It is a statement about representatives in two obstruction groups, not about the image of restriction from global objects.

\subsection{An initially free parameter can become an obstruction}
Consider the chain of views \(0\to1\to2\), with fiber \(\Z^2\) at each view. Write the maps in the direction in which data are restricted:
\[
 \rho_{0,1}(a,b)=(a,0),\qquad
 \rho_{1,2}(a,b)=(a,2b),\qquad
 \rho_{0,2}(a,b)=(a,0).
\]
The composite identity \(\rho_{0,1}\rho_{1,2}=\rho_{0,2}\) is literal. A global family is exactly
\begin{equation}\label{eq:branch-family}
 x_0=(k,0),\qquad x_1=(k,2n),\qquad x_2=(k,n),
 \qquad k,n\in\Z.
\end{equation}
Indeed the first coordinate must agree at every view. The second coordinate at view zero is forced to vanish, and the second coordinate at view one must be twice the second coordinate at view two. These conditions also suffice.

First prescribe only \(U=\{0\}\), with seed \((k,0)\). Every integer \(n\) in \eqref{eq:branch-family} gives a completion. Now prescribe \(W=\{0,1\}\), with values \((k,0)\) and \((k,m)\). This is a coherent partial section for every integer \(m\), because \(\rho_{0,1}\) forgets its second coordinate. Global extendability is a stronger requirement: \(m\) must be even.

\begin{proposition}[A coherent branch can be a dead end]\label{KS60}
Every initial seed \((k,0)\) in this model has a completion set in bijection with \(\Z\). A fixed intermediate value \((k,m)\) extends precisely when \(m\in2\Z\), and then has a unique completion. In particular the nonzero initial seed \((7,0)\) extends, while the coherent intermediate choice \((7,1)\) does not.
\end{proposition}
\begin{proof}
Equation~\eqref{eq:branch-family} supplies a completion for each \(n\), and the second coordinate of \(x_2\) recovers the parameter. Fixing the intermediate value imposes \(2n=m\). This equation has no integer solution for odd \(m\), and exactly one for even \(m\). Taking \(k=7\) and \(m=1\) proves the final assertion.
\end{proof}
\source{KS60}
The invisible coherent changes for the initial prescription are
\[
 ((0,0),(0,2n),(0,n)),\qquad n\in\Z.
\]
Thus \(H_U\simeq\Z\). Requiring invisibility also at view one forces \(2n=0\), hence \(n=0\), so \(H_W=0\). More observation removes all ambiguity in the successful branches. It does not turn the unsuccessful branches into solutions.

\plate{19_branches}{Every displayed middle value is internally coherent with the initial seed seven. Only even middle values lift to the last view. The initial problem has many global answers, but choosing an odd middle value destroys extendability. The arrows show actual integer preimages, not a probabilistic search policy.}{fig:branches}

\subsection{Why the extension space is a dependent sum}
Let \(s\) be the initial partial section. The stage theorem gives
\[
 \Ext_U(s)\simeq
 \coprod_{\substack{t\in\Gamma_W(A)\\t|_U=s}}\Ext_W(t).
\]
In the present example the index \(t\) is the integer \(m\). The summand is empty for odd \(m\) and a singleton for even \(m\). Taking their disjoint union gives one integer's worth of completions again. No contradiction appears between many initial completions and unique completion after a successful intermediate choice.

This formula also explains a logical error in greedy reconstruction. The fact that some branch extends does not justify replacing the existential choice by an arbitrary coherent branch. A method that fixes an intermediate value must either retain the remaining extension obligation or supply a rule proved to preserve extendability. Coherence on the current patch alone provides no such guarantee.

For an increasing sequence of prescribed domains, the same warning applies repeatedly. The finite stage formula keeps each intermediate choice explicit. It does not by itself manufacture one compatible infinite sequence of choices. The finite-fiber theorem later in the book supplies an actual existence criterion for one important infinite setting.

\subsection{Changing the coefficients changes both questions}
The equations behind the example can be read over a general abelian coefficient group \(T\). With the initial first coordinate fixed, the possible completions are parametrized by \(n\in T\). Prescribing the middle value \(m\) gives the equation \(2n=m\). Direct substitution therefore gives
\[
 \text{feasible middle values}=2T,\qquad
 \text{ambiguity after prescribing the middle}=\ker(2:T\to T).
\]
This calculation specializes as follows.
\begin{center}
\begin{tabular}{@{}lll@{}}
\toprule
Coefficient group & Which middle values extend? & Completions of an extendable value\\
\midrule
\(\Z\) & Even integers & One\\
\(\Q\) & Every rational & One\\
\(\mathbb F_2\) & Only zero & Two\\
\(\mathbb F_3\) & Every element & One\\
\bottomrule
\end{tabular}
\end{center}
The integer instance is the actual formal example cited above. The other rows are worked coefficient calculations from the same two equations. In characteristic two, stronger prescription does not eliminate the hidden parameter because doubling annihilates it. Over the rationals it creates no parity obstruction because doubling is invertible.

A solver that changes the coefficient domain while leaving its user-facing description untouched changes the mathematical problem. The rational answer \(1/2\) is not an integral completion of the odd middle value one. Conversely, a failure over \(\Z\) is not evidence of a failure over \(\Q\). The coefficient structure belongs beside the views, restrictions, and preserved values in the problem specification.

\begin{insight}
More information can reduce the set of answers without making the problem easier. The obstruction quotient records which new demands become impossible, while the ambiguity kernel records which freedom survives the demands that can be met.
\end{insight}
